% H9-T10: actual codimension-two algebraic representatives.
% New deduction in run JAC1964-20261121, by /root/h9_height.
% GPT-6 Astra Pro; effort setting not exposed.

\subsection*{H9-T10. The actual algebraic representative of the generic threefold}

\begin{theorem}[H9-T10]
Let $X/k$ be a smooth projective fourfold over an algebraically closed
field of arbitrary characteristic. Suppose that the Lefschetz lowering
operator of $X$ is algebraic, and let $Y/K$, $K=k(\mathbf P^1)$, be the
generic fibre of a genuine Lefschetz pencil of hyperplane sections of
$X$. Then the variable factor of the actual universal regular algebraic
representative $\operatorname{Ab}^2(Y)$ is polarized by the rational
incidence divisor $\Delta^*d/2$. Its canonical map to the variable
incidence Picard factor is an isogeny.

If $X$ is abelian and 032, Corollary 8.4, is accepted, the same assertions
hold for the whole primitive factor:
\[
 \operatorname{Ab}^2(Y)_{\mathrm{prim}}
          \longrightarrow P^2(Y)_{\mathrm{prim}}
       \quad\text{is an isogeny,}                           \tag{H9.32}
\]
and the normalized primitive incidence divisor on its source is a
rational polarization. If $X$ is a smooth complete intersection in
projective space, all of $H^3(Y)$ is variable and primitive, so the
assertion holds for the entire $\operatorname{Ab}^2(Y)$, without 032.
The assertion concerns the generic fibre and its displayed abelian
varieties after base change; it does not assert the same conclusion
for every special fibre.
\end{theorem}

\paragraph{Existence over the imperfect pencil field.}
The symbol $\operatorname{Ab}^2(Y)$ here denotes the actual universal
regular algebraic representative over $K$, even when $K=k(t)$ is
imperfect. Its existence is the arbitrary-field case of
Achter--Casalaina-Martin--Vial, \emph{A functorial approach to regular
homomorphisms}, \emph{Algebraic Geometry} 10 (2023), Theorem 2,
printed p.~89, proved as Theorem 6.1, printed p.~114. Take
$S=\Lambda=\operatorname{Spec}K$ in that theorem. The authors explicitly
state the new arbitrary-field codimension-two conclusion immediately
before Theorem 2. No perfectness assumption or descent from a geometric
object is needed for this existence assertion. We do not assert an
integral universal base-change theorem along purely inseparable
extensions.

\paragraph{The codimension-two realization input.}
We first prove the exact consequence of the known Bloch-map results
needed here. Let $V/F$ be smooth projective over any field, put
$R=\operatorname{Ab}^2(V)$, and let $B$ be an abelian factor of $R$ in
the rational category. Let $U\in\operatorname{CH}^2(B\times V)_{\mathbf Q}$
be a pointed family whose induced regular homomorphism is the inclusion
$\iota:B\to R$. Then, for every prime $\ell\ne\operatorname{char}F$,
its realization map is injective:
\[
 f_\ell:V_\ell B\lhook\joinrel\longrightarrow
                   H^3(V_{\overline F},\mathbf Q_\ell(2)).  \tag{H9.33}
\]

The sources are Achter--Casalaina-Martin--Vial, Ast\'erisque 455 (2025),
\S1.3 and Lemmas 1.3--1.4, printed pp.~6--7, and the proof of
Proposition 8.1, printed pp.~41--42. The first gives a miniversal cycle
over an arbitrary field for every surjective regular homomorphism.
Lemma 1.3 says that $b\mapsto[b]-[0]$ identifies the torsion of an
abelian variety with the torsion of its algebraically trivial zero-cycle
group. Lemma 1.4 consequently gives the compatibility on Tate modules
of a pointed family with its regular homomorphism. The proof of
Proposition 8.1 uses the injectivity of the codimension-two Bloch map
$V_\ell\lambda^2$, in every characteristic with $\ell$ invertible.

For completeness, these statements give (H9.33) directly. Clear a
positive denominator of $U$ and then tensor all Tate modules with
$\mathbf Q_\ell$. The family gives
\[
 V_\ell B\xrightarrow{V_\ell w_U}
 V_\ell A^2(V_{F^{\rm sep}})
 \xrightarrow{V_\ell\lambda^2}
 H^3(V_{\overline F},\mathbf Q_\ell(2)).
\]
Writing $\varphi$ for the universal regular homomorphism,
Lemma 1.4 gives
$(V_\ell\varphi)(V_\ell w_U)=V_\ell\iota$.
Since $\iota$ is an inclusion up to isogeny, $V_\ell w_U$ is
injective. The second arrow is injective by the cited Bloch-map
result. Their composite is the cohomological correspondence map
by the naturality used in Proposition 8.1, proving (H9.33).
The same input is also stated explicitly in the 2023 paper:
Theorem 8.1, printed p.~119, gives injectivity of the second
$\ell$-adic Bloch map over a separably closed field, and Proposition
8.2, printed p.~120, gives the miniversal realization injection at
primes away from its degree. The rational argument just given also
covers primes dividing that degree, since its positive integer is
invertible in $\mathbf Q_\ell$.
This proves injectivity for a chosen normalized miniversal realization;
it does not assert that $V_\ell\varphi$ is injective, or that
$V_\ell\lambda^2$ is surjective onto all of $H^3$.

\paragraph{Projectors and normalized families.}
Put $h=[Y]$. On $H^3(Y)$ the algebraic correspondence
\[
 \Pi_V=1-j^*(L_h:H^3(X)\longrightarrow H^5(X))^{-1}j_*
                                                               \tag{H9.34}
\]
is the self-adjoint variable projector. Choose the inverse supported
from degree $5$ to degree $3$, using the algebraic K\"unneth
projectors supplied by the Lefschetz standard theorem. On the other
degrees the subtracted correspondence is zero, so (H9.34) is still
an idempotent on the whole cohomology. Proposition 8.1 applied to
its idempotence defect gives an actual idempotent on
$R=\operatorname{Ab}^2(Y)$ in the isogeny category. Denote its image
by $B_V$.

For abelian $X$, let $e$ be the primitive projector of $H^3(X)$ onto
$\ker L_h^2$, and put
\[
 \Pi_F=j^*eL_h^{-1}j_*.
\]
Exactly as in H9-T4, this is an algebraic self-adjoint idempotent,
orthogonal to $\Pi_V$. It gives a factor $B_F$ of $R$. Their sum
acts on $H^3(Y)$ as the primitive projector. This defines the factor
$B_{\rm prim}$ in (H9.32), with
$B_{\rm prim}\sim B_F\times B_V$. The action on the algebraic
representative depends only on the degree-three action: to verify
this precise point, take a normalized miniversal family on $R\times Y$
and apply the parameter degree-one projector constructed below. If a
correspondence acts as zero on $H^3(Y)$, its composite with this family
has zero total cohomology class; Proposition 8.1 with source $R$ and
source codimension $\dim R$ makes the induced homomorphism $R\to R$
zero. Thus no off-middle action of the chosen representatives changes
the factor just defined.

Pull a miniversal family back to any of these factors, divide by its
positive degree, subtract its origin fibre, and apply the relevant
projector on $Y$. Its regular family map remains the factor inclusion.
As in JAC-T4.4, a rational polynomial $Q$ with $Q(2)=1$ and
$Q(2^a)=0$ for $0\le a\le2\dim B$, $a\ne1$, applied to
$([2]\times1)^*$, preserves this family map and leaves its cohomology
entirely in $H^1(B)\otimes H^3(Y)$. It has the required fixed or
variable image. This construction is valid in characteristic two:
the relevant cohomology uses $\ell\ne2$, the eigenvalues are $2^a$,
and the projector and family are rational correspondences.

\paragraph{The variable factor.}
If $B_V=0$, its assertion is empty. Otherwise (H9.33) gives an injection
\[
 V_\ell B_V\lhook\joinrel\longrightarrow E_\ell(2),
 \qquad E_\ell=\ker(j_*:H^3(Y)\to H^5(X)),                  \tag{H9.35}
\]
for every invertible prime. In particular $E_\ell\ne0$. SGA 7 II,
XVIII, 6.6 and 6.7.2, printed pp.~323--327, give spanning and conjugacy
of vanishing cycles and absolute irreducibility after every finite-index
monodromy restriction. Their parity hypotheses hold because $\dim X=4$;
the characteristic-two exception with odd ambient dimension is absent.
As in H9-T6a, all local vanishing cycles are nonzero and (H9.35) is an
isomorphism, even after the finite extension needed to define the family.

The cup product on $E_\ell$ is nondegenerate, because it is the
orthogonal complement of the nondegenerate ambient summand. Hence the
self-incidence divisor of the normalized family has a nondegenerate
alternating realization. Its associated symmetric homomorphism is
therefore an isogeny. This establishes the nondegeneracy here directly;
it was not borrowed from the incidence quotient.

The remaining sign argument is precisely H9-T6, now applicable to
$B_V$ and this family. To make clear which hypotheses have been
supplied, absolute irreducibility gives
$\operatorname{End}^0(B_{V,\overline K})=\mathbf Q$. If $L_0$ is a
genuine polarization, write
$\ell_V=\tfrac12\Delta^*d_V=cL_0$, $c\in\mathbf Q^*$.
The incidence Koszul sign gives
$E_{\ell_V}=-f_\ell^*Q_\ell=cE_0$. At one quadratic singular fibre,
after a common semistable extension of ramification degree $e>0$,
the comparison with the integral positive toric monodromy pairing of
H9-L2 gives
\[
 a_\ell^2=\frac{q_0}{e c}\quad(\ell\ne\operatorname{char}k),
 \qquad q_0\in\mathbf Q_{>0}
\]
with the same rational $q_0$ for every prime. H9-L1 forces $c>0$.
Thus $\ell_V$ is a rational polarization. All ingredients in this
argument are properties of the normalized family, its injective
realization and the Lefschetz monodromy; none requires that $B_V$
already be identified with an incidence Picard variety.

\paragraph{The fixed primitive factor for an abelian fourfold.}
Work over $\overline K$ to test ampleness. Let $C_0\to B_F$ be the
normalization of a nonconstant curve, pointed at the origin after
translation, and let $\xi$ be the class of its induced family.
The induced homomorphism $J(C_0)\to B_F$ is nonzero, so its rational
Tate-module map is nonzero by faithfulness. Composing with the
injection (H9.33) for $B_F$ remains nonzero. By naturality this
composite is the realization of the pulled-back family. Thus
$\xi\ne0$. This argument uses the base change of the existing
realization injection; it requires no claim that a universal regular
representative commutes integrally with purely inseparable base change.
The class lies in $H^1(C_0)\otimes F$, with $F$ as in H9-T4.

H9-T4 transports it, by the algebraic Lefschetz inverse on $X$ and the
Abel map of $C_0$, to a nonzero primitive codimension-two algebraic
class on the abelian variety $J(C_0)\times X$. Equation (H9.9) and
032 give
\[
 \deg_{C_0}\ell_F
       =\tfrac12\int_{C_0\times Y}\xi^2>0.
\]
Translation does not alter this degree, because it acts trivially on
the N\'eron--Severi group. Thus $\ell_F$ is nef and has positive degree
on every curve. If its associated symmetric homomorphism had a
positive-dimensional connected reduced kernel, the restriction to
that abelian subvariety would be algebraically trivial and a curve
inside it would have degree zero. This is impossible. The nef and
nondegenerate ampleness criterion proves that $\ell_F$ is a rational
polarization, exactly as in JAC-T4.4.

The fixed and variable cohomological subspaces are orthogonal. The
mixed incidence divisor between the corresponding families therefore
has zero mixed cohomology class. Faithfulness of the Tate module for
homomorphisms of abelian varieties makes the corresponding homomorphism
$B_F\to\widehat B_V$ zero. Thus their incidence forms are orthogonal
already in the rational N\'eron--Severi group. The product of the two
rational polarizations descends through the isogeny to polarize
$B_{\rm prim}$.

\paragraph{Comparison with the incidence quotient and the complete-intersection case.}
The incidence quotient map is regular, so universality gives a
surjection $R\to P^2(Y)$, compatible with algebraic correspondences.
Its variable and primitive restrictions are therefore surjective.
As in (JAC.T44.4), Saito's incidence identities give
\[
 d_B=(q\times q)^*d_{P,B}
       \quad\text{in }\operatorname{NS}(B\times B)_{\mathbf Q}.
\]
This can also be checked by replacing the normalized miniversal family
by the pullback of a normalized Poincar\'e family one factor at a time:
their difference has zero incidence map, so testing it against either
family kills the mixed divisor component. Since $\Delta^*d_B/2$ is
ample, the kernel of $q$ cannot contain a curve. Its kernel is finite,
which proves the isogeny claims, allowing a finite inseparable kernel
in positive characteristic.

If $X$ is a smooth four-dimensional complete intersection, weak
Lefschetz gives $H^3(X)=H^7(X)=0$. Hence $E_\ell=H^3(Y)$ and
$H^5(Y)=0$, so this cohomology is wholly primitive. The variable
projector acts as the identity on the algebraic representative by the
degree-three detection proved above. Its lowering operator is algebraic
by the explicit construction in JAC-P08. The variable proof alone
therefore gives the full assertion without 032.\hfill$\square$

\paragraph{Scope and status.}
This is a deduction written in this run from the cited codimension-two
Bloch injectivity, the source-checked SGA monodromy statements, and 032
where explicitly used. It supplies the actual algebraic representative
in the first higher-dimensional generic-pencil test. It asserts a
rational primitive polarization and an isogeny comparison, not an
integral principal polarization. No standard torsion assumption for
$\operatorname{Ab}^2$, no surjectivity of its realization onto all of
$H^3$, and no map on arbitrary homologically trivial cycles has been
assumed.

The primary 2023 existence and realization source is the publisher PDF
\url{https://api.algebraicgeometry.nl/Article/19846/2023-1-003.pdf}.
Theorem 2 (p.~89), Theorem 6.1 (p.~114), Theorem 8.1 (p.~119), and
Proposition 8.2 (p.~120) were directly read in this run. All additional
data used in the monodromy proof are algebraic and of finite type, so
a common finite field extension suffices to define them; the
finite-index irreducibility and common finite semistable extension
arguments are exactly those of H9-T6.
